% FILE: 01_research_question.tex
% Project: adversarial
% Thesis Role: adversarial-testing-and-hostile-assumption-verification

\section{Research Question}
\label{sec:adversarial-question}

\subsection{Problem Statement}
\label{sec:adversarial-question-problem}

Standard public process mining formats and classical discovery algorithms assume that event
logs are trustworthy by construction. XES (IEEE 14845) and OCEL 2.0 treat timestamps and
attributes as self-evident truths, operating under the assumption of a secure perimeter.
Neither standard mandates the manufacture of cryptographically unforgeable traces, requires
BLAKE3 hashing of event payloads, or links events causally in a hash chain
(\texttt{standards/ocel-adversarial-gravity.md}, Section 3).

This assumption is structurally untenable in high-stakes environments---M\&A due diligence,
regulatory audits, algorithmic governance---where adversaries with write access to an OCEL
SQLite database or JSON payload can arbitrarily forge records, alter attribute values at
specific timestamps, and inject phantom relationships into the event-object mapping table. The
modifications are indistinguishable from legitimate operations because there is no causal hash
chain linking the events. Classical discovery algorithms (Alpha Miner, Heuristics Miner,
Inductive Miner) have no mechanism to detect Sybil-injected case identifiers or
non-deterministic reality receipts, leading to invalid spaghetti models and incorrect fitness
and precision metrics (\texttt{sources/papers/adversarial-canon.md}).

The problem statement is therefore: existing process mining ingestion boundaries do not
enforce cryptographic provenance, and existing discovery algorithms do not assume adversarial
log conditions. Any system that relies on these mechanisms for board-admissible evidence is
architecturally negligent.

\subsection{Old-Regime Assumption Challenged}
\label{sec:adversarial-question-old-regime}

The old-regime assumption challenged by this corpus is that an event log is a faithful record
of a process by virtue of being emitted by a known system. Under this assumption, conformance
checking is a matter of aligning a discovered model against a log that is presumed truthful.
The adversarial corpus challenges this assumption at every level.

At the format level, the PROV-O gap in OCEL 2.0 means that \texttt{wasGeneratedBy},
\texttt{used}, and \texttt{wasAttributedTo} are not cryptographically enforced: a ``User ID:
123'' is merely a string, not a cryptographic signature proving authorization
(\texttt{standards/ocel-adversarial-gravity.md}, Section 2.2). At the ingestion level, a raw
pandas DataFrame is mutable, carries no cryptographic envelope, and must be rejected at the
boundary---a requirement that no standard process mining framework enforces. At the algorithm
level, canonical discovery algorithms structurally fail under Sybil attack because they rely
on temporal proximity and case identifiers for correlation without provenance verification.

\subsection{Hypothesis}
\label{sec:adversarial-question-hypothesis}

The hypothesis of the adversarial corpus is: a zero-trust ingestion boundary implementing
HMAC-SHA256 signature verification, SHA-256 hash chain integrity checks, timestamp
monotonicity gates, and impossible-velocity detection can reject all ten canonical attack
scenarios while admitting conforming traces---and this rejection can be verified through
executable Python \texttt{assert} statements that constitute falsifiable tests under the Van
der Aalst Constitution.

The secondary hypothesis is that the Chatman Constraint---strict injectivity $C(S_a) \neq
C(S_b)$ for all $S_a \neq S_b$ in the witness lattice---is a necessary structural law for
preventing witness lattice forgery in \texttt{wasm4pm-compat}, and that Witness Lattice
Monotonicity ($W_{\text{new}} = W_{\text{old}} \sqcup w_{\text{step}}$) is a sufficient
condition for detecting non-monotonic execution flow and triggering self-halt
(\texttt{sources/wasm4pm-compat/adversarial-type-law.md}, Laws I--III).

\subsection{Success Criteria}
\label{sec:adversarial-question-success}

Success is defined by the ADVERSARIAL\_RECEIPT in
\texttt{receipts/RECEIPT\_REGISTRY.md}: the adversarial corpus must contain a minimum of
three artifacts covering raw laundering, metric fabrication, and conformance theater
challenges. The corpus exceeds this minimum (actual: three core files plus extensive
supporting assets). Success also requires that the simulation engine's ten executable
scenarios produce the expected ACCEPTED/REJECTED verdicts verifiable by \texttt{assert}
statements, and that the standards audit script exits 0 on pass and 1 on fail. The ALIVE
OMEGA checkpoint
(\texttt{checkpoints/PROCESS\_INTELLIGENCE\_ADVERSARIAL\_V30.1.1\_OMEGA.md}) certifies that
all attack vectors, recursive logic bombs, and ontological paradoxes have been absorbed and
neutralized.
